\documentclass[a4paper,12pt]{article}

\usepackage[french]{babel}
\usepackage[T1]{fontenc}
\usepackage[utf8]{inputenc}
\usepackage{amsmath, amssymb}
\usepackage{geometry}
\usepackage{array}
\usepackage{tikz}
\usepackage{pgfplots}
\pgfplotsset{compat=1.18}

\geometry{top=2cm, bottom=2cm, left=2cm, right=2cm}

\begin{document}

\begin{center}\Large\textbf{CORRIGÉ DU SUJET 4}\end{center}

\section*{PREMIÈRE PARTIE : AUTOMATISMES -- QCM}

\subsection*{Question 1}
\[
\frac{11}{12} - \frac{12}{13} = \frac{11 \times 13 - 12^2}{12 \times 13} = \frac{143 - 144}{12 \times 13} = -\frac{1}{12 \times 13} < 0
\]
Donc $\dfrac{11}{12} < \dfrac{12}{13}$.

\boxed{\text{Réponse A}}

\emph{Astuces : pour comparer deux nombres, on peut étudier le signe de leur différence.}

\subsection*{Question 2}
1 m = 100 cm donc 1 m$^2$ = 100 $\times$ 100 = 10 000 cm$^2$.
Ainsi 2 m$^2$ = 20 000 cm$^2$.

\boxed{\text{Réponse D}}

\subsection*{Question 3}
$20 \times 0,3 = 2 \times 10 \times 0,3 = 2 \times 3 = 6$.

\boxed{\text{Réponse A}}

\subsection*{Question 4}
Le coefficient multiplicatif associé à une hausse de $30\%$ est $1 + \dfrac{30}{100} = 1,3$.
De plus : $12 \times 1,3 = 12 \times 1 + 12 \times 0,3 = 12 + 3,6 = 15,6$.

\boxed{\text{Réponse C}}

\subsection*{Question 5}
Pour déterminer graphiquement les antécédents de 1, on trace la droite horizontale $y = 1$. Elle coupe la courbe aux points d'abscisses $-3$ et $3$.

\boxed{\text{Réponse C}}

\emph{Méthode : pour déterminer graphiquement les antécédents d'un réel $\alpha$ par une fonction $f$, on trace la droite d'équation $y = \alpha$, on repère les éventuels points d'intersection avec la courbe et on lit les abscisses de ces points.}

\subsection*{Question 6}
Une probabilité est un nombre compris entre 0 et 1. $\dfrac{20}{19} \approx 1,05 > 1$ donc ce nombre ne peut pas être une probabilité.

\boxed{\text{Réponse B}}

\subsection*{Question 7}
Le diagramme montre les effectifs par classe d'âge :
- [0;20[ : 5 adhérents
- [20;40[ : 15 adhérents
- [40;60[ : 20 adhérents
- [60;80[ : 15 adhérents
Total : $5 + 15 + 20 + 15 = 55$ adhérents.

On ne peut pas calculer l'âge moyen exact, mais 27 ans est une valeur possible. La proposition A n'est donc pas assurément fausse.

L'effectif de la classe [60;80[ est 15, donc certains adhérents ont plus de 70 ans. La proposition B est vraie.

La moitié des adhérents correspond à $55 \div 2 = 27,5$. L'effectif cumulé jusqu'à [20;40[ est $5 + 15 = 20$, donc la médiane se situe dans la classe [40;60[. La proposition C « La moitié des adhérents a entre 20 et 40 ans » est fausse car seulement 15 adhérents (sur 55) sont dans cette classe.

\boxed{\text{Réponse C}}

\subsection*{Question 8}
Pour tout réel $a$ :
\[
(a+1)^2 - 3(a+1) = (a+1)\big[(a+1) - 3\big] = (a+1)(a-2)
\]

\boxed{\text{Réponse C}}

\subsection*{Question 9}
Le coefficient multiplicatif associé à une baisse de $20\%$ est $1 - \dfrac{20}{100} = 0,8$.
Le coefficient $c'$ associé à l'évolution réciproque est donc $c' = \dfrac{1}{0,8} = \dfrac{10}{8} = \dfrac{5}{4} = 1,25$.
Le pourcentage d'évolution réciproque est $t' = (1,25 - 1) \times 100 = 25\%$.

\boxed{\text{Réponse B}}

\subsection*{Question 10}
Une fonction affine non constante est de la forme $f(x) = ax + b$ avec $a \neq 0$. Sa représentation graphique est une droite non horizontale et non verticale, c'est-à-dire une droite oblique.

\boxed{\text{Réponse C}}

\subsection*{Question 11}
L'événement « obtenir un nombre inférieur ou égal à 4 » est réalisé par les issues 1, 2, 3 et 4.
Sa probabilité est la somme des probabilités élémentaires :
\[
0,1 + 0,1 + 0,25 + 0,3 = 0,75
\]

\boxed{\text{Réponse D}}

\newpage

\section*{DEUXIÈME PARTIE}

\subsection*{Exercice 1}

\begin{enumerate}
\item \textbf{Volume au printemps 2027 :}
\[
u_1 = 0,8 \times u_0 + 40 = 0,8 \times 1000 + 40 = 800 + 40 = 840 \text{ m}^3
\]

\item \textbf{Justification de la relation :}
\begin{itemize}
\item Perte de $20\%$ : on multiplie par $0,8$.
\item On ajoute ensuite $40$ m$^3$.
\item Donc $u_{n+1} = 0,8 u_n + 40$.
\end{itemize}

\item \textbf{La suite n'est pas arithmétique :}
\[
u_1 - u_0 = 840 - 1000 = -160
\]
\[
u_2 = 0,8 \times 840 + 40 = 672 + 40 = 712
\]
\[
u_2 - u_1 = 712 - 840 = -128 \neq -160
\]
La différence n'est pas constante, donc $(u_n)$ n'est pas arithmétique.

\item \textbf{Étude de la suite auxiliaire :}
\begin{enumerate}
\item $v_n = u_n - 200$
\[
v_{n+1} = u_{n+1} - 200 = (0,8 u_n + 40) - 200 = 0,8 u_n - 160
\]
\[
v_{n+1} = 0,8(u_n - 200) = 0,8 v_n
\]
Donc $(v_n)$ est géométrique de raison $q = 0,8$.

\item $v_0 = u_0 - 200 = 1000 - 200 = 800$
\[
v_n = v_0 \times q^n = 800 \times (0,8)^n
\]

\item $u_n = v_n + 200 = 800 \times 0,8^n + 200$
\end{enumerate}

\item \textbf{Analyse des valeurs :}
\begin{enumerate}
\item D'après le tableau, lorsque $n$ augmente, $u_n$ diminue et se rapproche de $200$. On peut conjecturer que la suite $(u_n)$ converge vers $200$ lorsque $n$ tend vers $+\infty$.

\item Le volume d'eau dans le bassin se stabilise autour de $200$ m$^3$ au bout de plusieurs années. Le bassin atteint un équilibre où les pertes sont compensées par l'apport annuel.
\end{enumerate}
\end{enumerate}

\newpage

\subsection*{Exercice 2}

\begin{enumerate}
\item \textbf{Calcul de la dérivée :}
\[
f(x) = \frac{3x - 3}{2x + 1}
\]
Forme $\dfrac{u}{v}$ avec $u(x) = 3x - 3$, $u'(x) = 3$, $v(x) = 2x + 1$, $v'(x) = 2$.
\[
f'(x) = \frac{u'v - uv'}{v^2} = \frac{3(2x+1) - (3x-3) \times 2}{(2x+1)^2}
\]
\[
f'(x) = \frac{6x + 3 - (6x - 6)}{(2x+1)^2} = \frac{6x + 3 - 6x + 6}{(2x+1)^2} = \frac{9}{(2x+1)^2}
\]

\item \textbf{Variations de $f$ :}
\[
f'(x) = \frac{9}{(2x+1)^2} > 0 \text{ pour tout } x \neq -\frac{1}{2}
\]
Donc $f$ est strictement croissante sur $]-\infty ; -\frac{1}{2}[$ et sur $]-\frac{1}{2} ; +\infty[$.

\item \textbf{Tangentes de coefficient directeur 1 :}
\begin{enumerate}
\item On résout $f'(x) = 1$ :
\[
\frac{9}{(2x+1)^2} = 1 \iff (2x+1)^2 = 9
\]
\[
2x+1 = 3 \quad \text{ou} \quad 2x+1 = -3
\]
\[
2x = 2 \implies x = 1 \quad \text{ou} \quad 2x = -4 \implies x = -2
\]
Deux solutions distinctes, donc exactement deux tangentes.

\item \textbf{Équations des tangentes :}
Formule : $y = f'(x_0)(x - x_0) + f(x_0)$

Pour $x_0 = 1$ :
\[
f(1) = \frac{3 \times 1 - 3}{2 \times 1 + 1} = \frac{0}{3} = 0
\]
\[
y = 1 \times (x - 1) + 0 = x - 1
\]

Pour $x_0 = -2$ :
\[
f(-2) = \frac{3 \times (-2) - 3}{2 \times (-2) + 1} = \frac{-6 - 3}{-4 + 1} = \frac{-9}{-3} = 3
\]
\[
y = 1 \times (x + 2) + 3 = x + 5
\]

Les deux tangentes ont pour équations : $y = x - 1$ et $y = x + 5$.
\end{enumerate}

\item \textbf{Graphique :}

\begin{center}
\begin{tikzpicture}[scale=0.6]
    \begin{axis}[
        axis lines = middle,
        xlabel = $x$,
        ylabel = $y$,
        xmin = -5, xmax = 5,
        ymin = -5, ymax = 5,
        xtick = {-4,-3,-2,-1,0,1,2,3,4},
        ytick = {-4,-3,-2,-1,0,1,2,3,4},
        grid = both,
        grid style = {gray!30},
        width = 12cm,
        height = 12cm,
        samples = 100,
        smooth,
        clip = false
    ]
        % Asymptote verticale
        \addplot[dashed] coordinates {(-0.5,-5) (-0.5,5)};
        
        % Courbe de f
        \addplot[thick, domain=-5:-0.6] {(3*x-3)/(2*x+1)};
        \addplot[thick, domain=-0.4:5] {(3*x-3)/(2*x+1)};
        
        % Tangentes
        \addplot[blue, thick, domain=-5:5] {x-1};
        \addplot[blue, thick, domain=-5:5] {x+5};
        
        % Points de tangence
        \addplot[only marks, mark=*, mark size=2, red] coordinates {(1,0) (-2,3)};
        
    \end{axis}
\end{tikzpicture}
\end{center}

\end{enumerate}

\vspace{1cm}

\textbf{Barème indicatif :}
\begin{itemize}
\item Automatismes : 6 points (0,5 pt par question)
\item Exercice 1 : 7 points
  \begin{itemize}
  \item 1) 1 pt
  \item 2) 1 pt
  \item 3) 1 pt
  \item 4a) 1 pt
  \item 4b) 1 pt
  \item 4c) 1 pt
  \item 5a) 0,5 pt
  \item 5b) 0,5 pt
  \end{itemize}
\item Exercice 2 : 7 points
  \begin{itemize}
  \item 1) 2 pts
  \item 2) 1,5 pts
  \item 3a) 1,5 pts
  \item 3b) 2 pts
  \end{itemize}
\end{itemize}

\end{document}